\begin{titlepage}
\thispagestyle{empty}
\begin{center}
\vspace*{1.4cm}
{\Large\bfseries [TÊN TRƯỜNG / KHOA — CẦN XÁC NHẬN]}\\[1.5cm]
{\Large\bfseries BÁO CÁO TỔNG KẾT BTL03}\\[0.6cm]
{\LARGE\bfseries CHUYỂN ĐỔI MÃ NGUỒN DEEP LEARNING\\ TỪ KERAS SANG PYTORCH}\\[0.8cm]
Snapshot: T0212\\
Commit nguồn: \texttt{33dcd74599e86404bbc8128c8057afba82810093}\\
Ngày chốt evidence: 2026-11-23

\vfill
\begin{tabularx}{\textwidth}{c p{0.26\textwidth} p{0.24\textwidth} X}
\toprule
STT & Thành viên / alias & Họ tên / MSSV & Vai trò / task evidence \\
\midrule
1 & Duy-DS & [CẦN NHÓM XÁC NHẬN] & T0101/T0201/T0202 reports; không suy ra tỷ lệ \\
2 & nguyenthanhdat180905-cloud & [CẦN NHÓM XÁC NHẬN] & Nhóm 2 theo task reports \\
3 & nguyen-cong-tri & [CẦN NHÓM XÁC NHẬN] & Nhóm 3 theo task reports \\
4 & TonOfWeight072501 & [CẦN NHÓM XÁC NHẬN] & Utilities theo task reports \\
5 & CHIVY2005 & [CẦN NHÓM XÁC NHẬN] & Nhóm 5 theo task reports \\
\midrule
\multicolumn{2}{l}{Giảng viên/lớp: [CẦN XÁC NHẬN]} & \multicolumn{2}{l}{Tỷ lệ đóng góp: [NHÓM XÁC NHẬN; tổng 100\%]} \\
\bottomrule
\end{tabularx}
\vfill
[ĐỊA ĐIỂM], 2026
\end{center}
\end{titlepage}

\chapter*{Tóm tắt}
\addcontentsline{toc}{chapter}{Tóm tắt}
BTL03 tái cấu trúc mười cấu hình thuộc năm nhóm Keras sang PyTorch. Tại T0212, report ghi parameter-count audit và 10/10 synthetic one-epoch smoke cases. Đây là kiểm tra cấu trúc/wiring, không phải training metric trên dữ liệu thật. Logger bị BLOCKED do protobuf/TensorBoard; full suite theo shard ghi 78 passed và 4 failed trên 82 test. T0212 chưa đủ điều kiện nghiệm thu/merge. T0201/T0202 là kết quả riêng, không đại diện full training cho mười cấu hình.

\noindent\textbf{Từ khóa:} Keras; PyTorch; parameter audit; reproducibility; synthetic smoke-test; provenance.

\chapter*{Danh mục viết tắt}
\addcontentsline{toc}{chapter}{Danh mục viết tắt}
\begin{tabularx}{\textwidth}{lY}
\toprule Viết tắt & Ý nghĩa \\
\midrule
CNN & Convolutional Neural Network \\
FCN & Score-map model theo legacy; không mặc nhiên là mAP detector \\
NCHW/NHWC & Batch–channel–height–width / batch–height–width–channel \\
OOP & Lập trình hướng đối tượng \\
PR & Pull Request / nhánh tích hợp được audit \\
RQ & Câu hỏi nghiên cứu \\
\bottomrule
\end{tabularx}

\chapter*{Câu hỏi nghiên cứu}
\addcontentsline{toc}{chapter}{Câu hỏi nghiên cứu}
Nguồn Project Charter Markdown đúng tên và RQ nguyên văn chưa được xác minh tại snapshot. Không dựng RQ từ README. Chỉ điền ma trận khi có charter chính thức:
\begin{longtable}{p{0.12\textwidth}p{0.43\textwidth}p{0.28\textwidth}p{0.14\textwidth}}
\toprule RQ ID & Câu hỏi nguyên văn từ charter đã xác minh & Evidence cần & Trạng thái \\
\midrule
[RQ ID] & [Chờ xác minh] & [Run/report/test] & [Chưa kết luận nếu thiếu run] \\
\bottomrule
\end{longtable}

\chapter*{Mở đầu}
\addcontentsline{toc}{chapter}{Mở đầu}
\section*{Bối cảnh và mục tiêu}
BTL03 khảo sát chuyển đổi notebook/mô hình thực hành từ Keras sang PyTorch OOP. Báo cáo tập trung snapshot T0212, phân biệt implementation với thực nghiệm còn bị chặn.

\section*{Đối tượng, phạm vi và phương pháp}
Đối tượng gồm mười cấu hình trong năm nhóm; phạm vi dừng tại commit T0212 nêu trên. Phương pháp là đối chiếu task reports, notebook/CSV provenance, source/config, tests và Git history; không suy diễn metric thiếu nguồn.

\section*{Đóng góp đã xác minh và giới hạn}
Nêu deliverable có nguồn; giữ nguyên trạng thái BLOCKED, synthetic smoke, logger blocker và task chưa nghiệm thu.

\section*{Cấu trúc báo cáo}
Chương 1 tổng quan; chương 2 khác biệt framework; chương 3 kiến trúc/pipeline; chương 4 audit và trạng thái; chương 5 kết luận/giới hạn/hướng tiếp theo.